% Contact cube (Tutorial 2): background mesh and GIMP domains, in the style of
% column_setup.png. The cube is 0.8 x 0.8 x 0.8 m: 2 x 2 x 2 elements of size
% h = 0.4 m, each holding a 2 x 2 x 2 grid of GIMPs.
% Build: pdflatex cube_setup.tex
%        pdftocairo -png -transp -r 300 -singlefile cube_setup.pdf cube_setup
\documentclass[border=4pt]{standalone}
\usepackage{tikz}
\usetikzlibrary{arrows.meta}

\definecolor{gridgrey}{RGB}{185,185,185}
\definecolor{activeorange}{RGB}{245,130,32}
\definecolor{mpnavy}{RGB}{0,0,153}
\definecolor{gimpfill}{RGB}{179,179,226}
\definecolor{gimpside}{RGB}{158,158,212}
\definecolor{gimptop}{RGB}{204,204,238}
\definecolor{support}{RGB}{254,238,224}
\definecolor{supportside}{RGB}{250,224,201}
\definecolor{supporttop}{RGB}{255,246,238}
\definecolor{dimpurple}{RGB}{104,36,109}

% lengths in units of h
\newcommand{\gimpside}{0.475} % drawn GIMP side (true side 0.5)

\tikzset{
  mesh/.style={draw=gridgrey, line width=1pt},
  hidden/.style={draw=gridgrey, line width=0.6pt, dash pattern=on 2.5pt off 2pt},
  dimline/.style={draw=dimpurple, line width=1pt},
  dim/.style={dimline, {Latex[length=5pt]}-{Latex[length=5pt]}},
  gimpedge/.style={draw=mpnavy, line width=0.5pt, line join=round},
  activenode/.style={fill=activeorange, draw=black, line width=0.5pt},
}

\newcommand{\activenode}[1]{\path[activenode] #1 circle (3pt);}
\newcommand{\mpoint}[1]{\fill[mpnavy] #1 circle (1.4pt);}

% the mesh lines inside the cube or on its back, left and bottom faces
\newcommand{\hiddenmesh}{%
  \foreach \a in {0,1} {
    \foreach \b in {1,2} {
      \draw[hidden] (0,\b,\a) -- (2,\b,\a);   % along x
      \draw[hidden] (\a,\b,0) -- (\a,\b,2);}  % along z
    \foreach \b in {0,1} {
      \draw[hidden] (\a,0,\b) -- (\a,2,\b);}} % along y
  \foreach \x in {0,1} {\foreach \y in {1,2} {\foreach \z in {0,1} {
    \path[activenode, opacity=0.45] (\x,\y,\z) circle (3pt);}}}}
% the mesh lines on the front, right and top faces
\newcommand{\visiblemesh}{%
  \foreach \a in {0,1,2} {
    \draw[mesh] (0,0,\a) -- (2,0,\a) -- (2,2,\a);   % front and right, across
    \draw[mesh] (\a,0,0) -- (\a,0,2) -- (\a,2,2);   % front and top, up
    \draw[mesh] (2,\a,0) -- (2,\a,2) -- (0,\a,2);}  % right and top, up
  \foreach \a in {0,1,2} {\foreach \b in {0,1,2} {
    \activenode{(\a,0,\b)}}}                        % front face
  \foreach \a in {1,2} {\foreach \b in {0,1,2} {
    \activenode{(2,\a,\b)}}}                        % right face
  \foreach \a in {0,1} {\foreach \b in {1,2} {
    \activenode{(\a,\b,2)}}}}                       % top face
\newcommand{\dimensions}{%
  \draw[dim] (-0.4,0,0) -- node[left=3pt, font=\sffamily] {0.8 m} (-0.4,0,2);
  \draw[dim] (0,0,-0.35) -- node[below=1pt] {$h$} (1,0,-0.35);
  \draw[dim] (2.35,0,0) -- node[below right=0pt, font=\sffamily] {0.8 m} (2.35,2,0);}

% one GIMP domain, lower corner (#1,#2,#3), side #4: the material point at the
% centre, then the three visible faces
\newcommand{\gimp}[4]{%
  \mpoint{(#1+#4/2,#2+#4/2,#3+#4/2)}
  \filldraw[gimpedge, fill=gimpfill, fill opacity=0.4]
    (#1,#2,#3) -- (#1+#4,#2,#3) -- (#1+#4,#2,#3+#4) -- (#1,#2,#3+#4) -- cycle;
  \filldraw[gimpedge, fill=gimpside, fill opacity=0.4]
    (#1+#4,#2,#3) -- (#1+#4,#2+#4,#3) -- (#1+#4,#2+#4,#3+#4) -- (#1+#4,#2,#3+#4) -- cycle;
  \filldraw[gimpedge, fill=gimptop, fill opacity=0.4]
    (#1,#2,#3+#4) -- (#1+#4,#2,#3+#4) -- (#1+#4,#2+#4,#3+#4) -- (#1,#2+#4,#3+#4) -- cycle;}

\begin{document}
% legend glyphs, drawn outside the main picture so they use plain axes
\newsavebox{\legmp}\sbox{\legmp}{\tikz[baseline=-0.6ex]\fill[mpnavy] (0,0) circle (1.4pt);}
\newsavebox{\leggimp}\sbox{\leggimp}{\tikz[baseline=-0.6ex]\filldraw[gimpedge, fill=gimpfill, fill opacity=0.55] (0,-0.15) rectangle (0.3,0.15);}
\newsavebox{\legnode}\sbox{\legnode}{\tikz[baseline=-0.6ex]\path[activenode] (0,0) circle (3pt);}
\newsavebox{\legsupport}\sbox{\legsupport}{\tikz[baseline=-0.6ex]\fill[support] (0,-0.15) rectangle (0.3,0.15);}

\begin{tikzpicture}[x={(1.5cm,0cm)}, y={(0.75cm,0.6cm)}, z={(0cm,1.5cm)}]

  % (a) background mesh, with the elements shaded as the support
  \begin{scope}
    \fill[support] (0,0,0) -- (2,0,0) -- (2,0,2) -- (0,0,2) -- cycle;
    \fill[supportside] (2,0,0) -- (2,2,0) -- (2,2,2) -- (2,0,2) -- cycle;
    \fill[supporttop] (0,0,2) -- (2,0,2) -- (2,2,2) -- (0,2,2) -- cycle;
    \hiddenmesh
    \visiblemesh
    \dimensions
    \node[font=\sffamily] at (1.25,0,-1.0) {(a) background mesh};
  \end{scope}

  % (b) a 2 x 2 x 2 grid of GIMP domains in each element, drawn back to front
  \begin{scope}[xshift=7cm]
    \hiddenmesh
    \pgfmathsetmacro{\gap}{(0.5-\gimpside)/2}
    \foreach \j in {3,2,1,0} {
      \foreach \i in {0,1,2,3} {
        \foreach \k in {0,1,2,3} {
          \gimp{0.5*\i+\gap}{0.5*\j+\gap}{0.5*\k+\gap}{\gimpside}}}}
    \visiblemesh
    \dimensions
    \node[font=\sffamily] at (1.25,0,-1.0) {(b) GIMP domains};
  \end{scope}

  % legend, centred under the two panels
  \node[font=\sffamily, anchor=north] at (5.6cm,-1.95cm) {%
    \usebox{\legmp}~material point\hspace{1.6em}%
    \usebox{\leggimp}~GIMP domain\hspace{1.6em}%
    \usebox{\legnode}~active node\hspace{1.6em}%
    \usebox{\legsupport}~support};

\end{tikzpicture}
\end{document}
